% =============================================================
\section{Cell models}\label{sec:cells}
% =============================================================

\subsection{Hodgkin--Huxley straight cable}\label{sec:cells-hh}

The verification studies of \cref{sec:results-verify} and the pilot
strength--duration sweep of \cref{sec:res-sd} use a uniform
Hodgkin--Huxley cable. Defaults reproduce the configuration set by
\texttt{experiment.make\_hh\_cable\_experiment}: $50$ compartments
along a $1.25\,$mm cable, fibre radius $10\,\mu$m, axial resistivity
$\rho_a = 100\,\Omega\cdot$cm, with the standard Jaxley HH channel set on
every compartment. A single electrode sits $50\,\mu$m above the first
compartment in a medium of conductivity $\sigma = 0.3\,$S/m, the typical
value for cortical grey matter.

\subsection{BBP L2/3 cells (cADpyr229, cNAC187)}\label{sec:cells-bbp}

For morphologically detailed simulations we adapt two L2/3 cell models
from the Blue Brain Project~\cite{markram2015reconstruction}: the
regular-spiking pyramidal cell
\texttt{L23\_PC\_cADpyr229} and the parvalbumin-expressing basket cell
\texttt{L23\_LBC\_cNAC187}. Channel kinetics are reused from
\texttt{jaxley\_mech.channels.l5pc}, which provides Jaxley-compatible
implementations of the BBP L5 channel set; conductances and passive
parameters are read from the BBP biophysics specification and applied per
section group (soma, axon, apical, basal).

Two implementation details affect reproducibility but do not bear on
the paper's results: the SWC conversion of the BBP soma contour, and
the calcium-channel insertion order in Jaxley. Both are documented in
\cref{app:bbp-implementation}.

A third detail is needed for NEURON parity. The \texttt{cADpyr229}
apical $I_h$ conductance follows a distance-dependent gradient
specified in BBP's biophysics file with a path-distance argument; we
apply it under a Euclidean-distance approximation to the soma centre,
which closes a $0.6\,$mV resting-potential offset that a uniform
somatic value otherwise leaves (closed-form gradient and
implementation in \cref{app:ih-gradient}). Long SWC branches are
subdivided at load time (\texttt{max\_branch\_len}$\,=\,100\,\mu$m) so
the compartment count tracks what NEURON's \texttt{geom\_nseg} routine
produces for the same morphology; we found this necessary for
spike-timing parity (\cref{sec:discussion}).

